\documentclass[a4paper]{article}
\usepackage[utf8]{inputenc}


%=-=-=-=-=-=-=-=-=-=-=-=-=-=-=-=-=-=-=-=-=-=-=-=-=-=-=-=-=-=-=-=-=-=-=-=-=-=-=-=-
% PREAMBLE
%=-=-=-=-=-=-=-=-=-=-=-=-=-=-=-=-=-=-=-=-=-=-=-=-=-=-=-=-=-=-=-=-=-=-=-=-=-=-=-=-

%%%%%%%%%%%%%%%%%%%%%%%%%%%%%%%%%%%%%%%%%%%%%%%%%%%%%%%%%%%%%%%%%%%%%
% Important styling notes
%%
% For now, to include img.jpg in img/path/to/img.jpg, just use:
% path/to/img.jpg - for details see style.tex
%=-=-=-=-=-=-=-=-=-=-=-=-=-=-=-=-=-=-=-=-=-=-=-=-=-=-=-=-=-=-=-=-=-=-=-=-=-=-=-=-
% Packages
%%
%\usepackage{fullpage} % Package to use full page
\usepackage[top=1in,bottom=1in,left=1in,right=0.7in,heightrounded]{geometry}

\usepackage{parskip}                    % Package to tweak paragraph skipping
\usepackage{amsmath}                    % standard
\usepackage{amssymb}                    % standard - Double R symbol etc.
\usepackage{hyperref}
\usepackage{amsthm}                     % standard - theorem, definition, etc.
\usepackage{multicol}                   % multiple columns for numbering
\usepackage{enumitem}                   % standard - enumerate styles
\usepackage[utf8]{inputenc}
\usepackage{scrextend}                  % indentation
\usepackage{graphicx}                   % standard - add figures
\usepackage{float}                      % standard - figure position, use [H] option
\usepackage{pifont}                     % symbols
\usepackage{gensymb}                    % degree symbol \degree
\usepackage{xcolor}                     % bg color
\hypersetup{
    colorlinks,
    linkcolor={black!50!black},
    citecolor={blue!50!black},
    urlcolor={blue!80!black}
}
\usepackage{framed}                     % bg color
\usepackage[T1]{fontenc}                % small caps
\usepackage{sectsty}                    % headings colour
\usepackage{mathtools}                  % Loads amsmath
\usepackage{amsthm,thmtools,xcolor}     % coloured theorem
\usepackage[toc,page]{appendix}         % reference to appendix
%\usepackage{titlesec}                   % change chapter, section, etc. formats
\usepackage{xifthen}                    % if, else
\usepackage{etoolbox}
% format numbering in theorem, lemma, etc. environment
\AtBeginEnvironment{theorem}{\setlist[enumerate, 1]{font=\upshape,  wide=0.5em, before=\leavevmode}}
\AtBeginEnvironment{lemma}{\setlist[enumerate, 1]{font=\upshape,  wide=0.5em, before=\leavevmode}}
\usepackage[letterspace=150]{microtype} % \textls{<letterspaced text>} % 0 <= letterspace <= 1000, 1000 = M space
\usepackage{letltxmacro}                % renew commands?
\usepackage{minted}                     % package to list code
    % otherwise minted goes off the page
    \setmintedinline{breaklines}
\usepackage{subfig}
\usepackage{eso-pic}                    % title page bg pic
\usepackage{varwidth}
\PassOptionsToPackage{svgnames}{xcolor}
\usepackage{fontawesome}                % \faQuestionCircle
\usepackage{marvosym}                   %\Pointinghand
\usepackage{mdframed}                   % easy outline frames
\usepackage[many]{tcolorbox}            % colour box for theorem styles
\usepackage{array,booktabs,calc} % table figs and text
\usepackage{comment}                    % \begin{comment}
\usepackage{fancyhdr}                   % page headings
\usepackage{mdframed}                   % boxes
\usepackage[backend=biber,sorting=none,style=ieee]{biblatex}
\usepackage{caption}
%%% caption options {
%\DeclareCaptionFont{white}{\color{white}}
\DeclareCaptionFormat{listing}{\colorbox{magenta!30!gray}{\parbox{\textwidth}{#1#2#3}}}
\captionsetup[lstlisting]{format=listing,labelfont={bf,small},textfont=small,skip=-1pt}
%%% }
\addbibresource{bibliography.bib}
\usepackage{url}
\usepackage{textcomp}
\usepackage[makeroom]{cancel}            % crossed symbols
\usepackage{algorithm}
\usepackage[noend]{algpseudocode}
\usepackage{tikz}
\usetikzlibrary{arrows.meta,positioning,quotes} % arrows and nodes in tikz
\usepackage{marginnote}
\usepackage{pgfplots}
\usepackage{pstricks-add,pst-slpe}  % for fancy tikz arrows
%\usepackage{titlesec}                   % title style
\usepackage{lmodern}                    % a font
\usepackage{titletoc} % Required for manipulating the table of contents
\usepackage{titlesec} % Allows customization of titles
\usepackage{fouriernc} % Use the New Century Schoolbook font
\usepackage{booktabs} % things in page margins
\usepackage{stmaryrd } % \varoast
\usepackage{listings} % code listings
\usepackage{longtable} % table across multiple pages
\usepackage{styles/nasm/lang}  % include custom language for NASM assembly.
\usepackage{styles/nasm/style} % include custom style for NASM assembly.



%% extra comments that I don't know where they belong:
% list of ding tags: http://willbenton.com/wb-images/pifont.pdf

%=-=-=-=-=-=-=-=-=-=-=-=-=-=-=-=-=-=-=-=-=-=-=-=-=-=-=-=-=-=-=-=-=-=-=-=-=-=-=-=-
% Colours for various things
%%


\definecolor{shadecolor}{rgb}{1.,0.933,0.96} % bg color, r,g,b <= 1
\definecolor{medium_blue}{RGB}{60,125,190}
\definecolor{dark_blue}{RGB}{25,60,85}
\definecolor{dark_red}{RGB}{77,16,16}
\definecolor{LightPink}{rgb}{0.92.,0.8,0.84} % bg color, r,g,b <= 1
\definecolor{LighterPink}{rgb}{1.,0.94,0.97} % bg color, r,g,b <= 1
\definecolor{LightestPink}{rgb}{1.,0.95,0.99} % bg color, r,g,b <= 1
\definecolor{DarkestPink}{rgb}{0.36, 0.0, 0.18}
\definecolor{DarkerPink}{rgb}{0.41, 0.0, 0.21}
\definecolor{DarkPink}{rgb}{0.55, 0.05, 0.37}
\definecolor{lightestestpink}{RGB}{255,248,252}
\definecolor{codegray}{rgb}{0.5,0.5,0.5}
\definecolor{codegrayblue}{rgb}{0.35,0.35,0.47}



%=-=-=-=-=-=-=-=-=-=-=-=-=-=-=-=-=-=-=-=-=-=-=-=-=-=-=-=-=-=-=-=-=-=-=-=-=-=-=-=-
% Define my own theorem styles
%%

% "base" styles
\declaretheoremstyle[
  headfont=\color{DarkPink}\bfseries,
  bodyfont=\itshape,
]{colored}

\declaretheoremstyle[
  headfont=\color{DarkPink}\bfseries,
  bodyfont=\normalfont,
]{colored_upright}

% theorems (corollaries, etc) themselves, inherit from my style above
% Usage:
% \begin{theorem} \end{theorem}, \begin{lemma} \end{lemma}, ...
\declaretheorem[
	numberwithin=section,
 	style=colored,
	name=\textsc{Theorem},
]{theorem}

\tcolorboxenvironment{theorem}{
  boxrule=0pt,
  boxsep=2pt,
  colback={magenta!25!white},
  colframe=DarkPink,
  enhanced jigsaw, 
  borderline west={2pt}{0pt}{DarkPink},
  sharp corners,
  before skip=5pt,
  after skip=5pt,
  breakable,
  right=0mm % for equations
}

\declaretheorem[
	numberwithin=section,
 	style=colored,
	name=\textsc{Corollary},
]{corollary}

\tcolorboxenvironment{corollary}{
  boxrule=0pt,
  boxsep=1pt,
  colback={magenta!10!white},
  colframe=DarkPink,
  enhanced jigsaw, 
  borderline west={2pt}{0pt}{DarkPink},
  sharp corners,
  before skip=5pt,
  after skip=5pt,
  breakable,
  right=0mm % for equations
}

\declaretheorem[
	numberwithin=section,
	style=colored,
	name=\textsc{Lemma},
]{lemma}

\tcolorboxenvironment{lemma}{
  boxrule=0pt,
  boxsep=1pt,
  colback={magenta!10!white},
  colframe=DarkPink,
  enhanced jigsaw, 
  borderline west={2pt}{0pt}{DarkPink},
  sharp corners,
  before skip=5pt,
  after skip=5pt,
  breakable,
  right=0mm % for equations
}

\declaretheorem[
	numberwithin=section,
	style=colored,
	name=\textsc{Definition},
]{definition}

\tcolorboxenvironment{definition}{
  boxrule=0pt,
  boxsep=1pt,
  colback={magenta!25!white},
  colframe=DarkPink,
  enhanced jigsaw, 
  borderline west={2pt}{0pt}{DarkPink},
  sharp corners,
  before skip=5pt,
  after skip=5pt,
  breakable,
  right=0mm % for equations
}

\declaretheorem[
	numberwithin=section,
  	style=colored,
  	name=\textsc{Example},
]{exmp}

\declaretheorem[
	numberwithin=section,
  	style=colored,
  	name=\textsc{Solution},
]{soln}

%%% code listings
\lstdefinestyle{code1}{
    backgroundcolor=\color{lightestestpink},   
    commentstyle=\color{codegrayblue},
    keywordstyle=\color{DarkerPink},
    numberstyle=\tiny\color{codegray},
    stringstyle=\color{black!40!cyan},
    basicstyle=\small\ttfamily,
    breakatwhitespace=false,
    breaklines=true,        
    captionpos=t,             
    keepspaces=true,        
    numbers=left,           
    numbersep=5pt,
    showspaces=false, 
    showstringspaces=false,
    showtabs=false,
    tabsize=4
}

\lstset{style=code1}

%=-=-=-=-=-=-=-=-=-=-=-=-=-=-=-=-=-=-=-=-=-=-=-=-=-=-=-=-=-=-=-=-=-=-=-=-=-=-=-=-
% Headers (size, font, colour)
%%




\makeatletter
\renewcommand{\@seccntformat}[1]{\llap{\textcolor{DarkestPink}{\csname the#1\endcsname}\hspace{1em}}}                    
\renewcommand{\section}{\@startsection{section}{1}{\z@}
{-4ex \@plus -1ex \@minus -.4ex}
{1ex \@plus.2ex }
{\normalfont\large\sffamily\bfseries\textcolor{DarkestPink}}}
\renewcommand{\subsection}{\@startsection {subsection}{2}{\z@}
{-3ex \@plus -0.1ex \@minus -.4ex}
{0.5ex \@plus.2ex }
{\normalfont\sffamily\bfseries\textcolor{DarkestPink}}}
\renewcommand{\subsubsection}{\@startsection {subsubsection}{3}{\z@}
{-2ex \@plus -0.1ex \@minus -.2ex}
{.2ex \@plus.2ex }
{\normalfont\small\sffamily\bfseries\textcolor{DarkestPink}}}                        


%=-=-=-=-=-=-=-=-=-=-=-=-=-=-=-=-=-=-=-=-=-=-=-=-=-=-=-=-=-=-=-=-=-=-=-=-=-=-=-=-
% Numberings, counters and spacings
%%
\numberwithin{equation}{section} % section number in eq/s
\setlength{\jot}{7pt} % spacing in split, gathered env/s



%% Custom examples
%% Output - Example 1,2,...
\newcounter{example}
\newenvironment{example}[1][]{\refstepcounter{example}\par\medskip
   \textbf{Example~\theexample. #1} \rmfamily}{\medskip}
%%%%%%%%%%%% End of unused %%%%%%%%%%%%



%=-=-=-=-=-=-=-=-=-=-=-=-=-=-=-=-=-=-=-=-=-=-=-=-=-=-=-=-=-=-=-=-=-=-=-=-=-=-=-=-
% Paths
%%
\graphicspath{ {./img/} } % figures' path - can look up files directly from there


%=-=-=-=-=-=-=-=-=-=-=-=-=-=-=-=-=-=-=-=-=-=-=-=-=-=-=-=-=-=-=-=-=-=-=-=-=-=-=-=-
% User defined macros (math mode)
%%


% Curly braces under text. Usage: \myunderbrace{upper}{lower}
\newcommand{\myunderbrace}[2]{\mathrlap{\underbrace{\phantom{#1}}_{#2}} #1}
\newcommand{\setR}{\mathbb{R}} % \ouble R
\newcommand{\setRn}{\mathbb{R}^n} %  double R^n
\newcommand{\setN}{\mathbb{N}} % double N
\newcommand{\setZ}{\mathbb{Z}} % double Z
\let\oldemptyset\emptyset
\let\emptyset\varnothing % nice - looking empty set symbol
\newcommand{\fancyN}{\mathcal{N}} % null space
\newcommand{\fancyR}{\mathcal{R}} % range

\newcommand{\bx}{\textbf{x}}
\newcommand{\by}{\textbf{y}}
\newcommand{\bb}{\textbf{b}}
\newcommand{\bA}{\textbf{A}}
\newcommand{\bB}{\textbf{B}}
\newcommand{\bI}{\textbf{I}}
% double bars as in norm
\newcommand{\norm}[1] {\lVert #1 \rVert} 
\newcommand{\trans}[1]{#1^{\top}}

\newcommand{\mean}[1]{\bar{#1}}
\newcommand{\var}{\sigma^2}

\newcommand{\partdevx}[1]{\frac{\partial #1}{\partial x}}
\newcommand{\partdevxx}[1]{\frac{\partial #1}{\partial x}}
\newcommand{\partdevxn}[1]{\frac{\partial^n #1}{\partial x^n}}
\newcommand{\partdevy}[1]{\frac{\partial #1}{\partial x}}
\newcommand{\partdevyy}[1]{\frac{\partial #1}{\partial y}}
\newcommand{\partdevyn}[1]{\frac{\partial^n #1}{\partial y^n}}

% text above = symbol
\newcommand{\overeq}[1]{\ensuremath{\stackrel{#1}=}} 
\newcommand{\greatersmaller}{%
  \mathrel{\ooalign{\raisebox{.6ex}{$>$}\cr\raisebox{-.6ex}{$<$}}}
} % greater and smaller symbols on top of each other, same line

%=-=-=-=-=-=-=-=-=-=-=-=-=-=-=-=-=-=-=-=-=-=-=-=-=-=-=-=-=-=-=-=-=-=-=-=-=-=-=-=-
% User defined macros (non math)

\newcommand{\qedblack}{$\hfill\blacksquare$} % black square end of line
\newcommand{\qedwhite}{\hfill \ensuremath{\Box}} % white square end of line
\newcommand{\hquad}{\hskip0.5em\relax}% half quad space
%\newcommand{\TODO}{\textcolor{red}{\bf TODO!}\;}

\newcommand{\TODO}[1][]{%
    \ifthenelse{\equal{#1}{}}{\textcolor{red}{\bf TODO!}\;}{\textcolor{red}{\textbf {TODO:} #1}\; }%
}
\newcommand{\B}[1]{\textbf{\textup{#1}}} % bold and upright
\renewcommand{\labelitemi}{\scriptsize$\textcolor{DarkPink}{\blacksquare}$} % itemize - squares instead of bullets
\newcommand{\emphasis}[1]{\textls{#1}}

\LetLtxMacro{\originaleqref}{\eqref}
\renewcommand{\eqref}{Eq.~\originaleqref}
\renewcommand*{\eqref}[1]{Eq.~\originaleqref{#1}}





% background images
%%%%%%%
\newcommand\BackgroundPic{%
\put(0,0){%
\parbox[b][\paperheight]{\paperwidth}{%
\vfill
%\centering
\includegraphics[width=0.125\paperwidth,height=\paperheight,%
]{img/background_02.png}% use ,keepaspectratio
\vfill
}}}
%%%%%%%
% end of background image
%%%%%%%%%%%%%% my own frame
\newmdenv[topline=false,bottomline=false]{leftrightbox}
%%%%%%%%%%%%% end
%%%%%%%%%%%%% my own comment
\newcommand{\mycomment}[1]{\begin{leftrightbox}\Pointinghand~\textbf{Comment:}~#1 \end{leftrightbox}}
%%%%%%%%%%%%% end
% my custom note https://tex.stackexchange.com/questions/301993/create-custom-note-environment-with-tcolorbox
\newmdenv[
    topline=false,
    bottomline=false,
    rightline=false,
    innerrightmargin=0pt
]{siderule}
\newenvironment{mynote}%
    {\begin{siderule}\textbf{\Pointinghand~Note:}}
    {\end{siderule}}
%%%%%%%%%%%%% my own box
\newcommand{\boxone}[1]{\begin{tcolorbox}[colback = LighterPink,colframe=LightPink]
#1
\end{tcolorbox}}
%%%%%%%%%%%%% end

\let\oldemptyset\emptyset
\let\emptyset\varnothing
%algorithmic
\algdef{SE}[DOWHILE]{Do}{doWhile}{\algorithmicdo}[1]{\algorithmicwhile\ #1}%






\begin{document}
%=-=-=-=-=-=-=-=-=-=-=-=-=-=-=-=-=-=-=-=-=-=-=-=-=-=-=-=-=-=-=-=-=-=-=-=-=-=-=-=-
% GLOBAL STYLES (DOCUMENT SCOPE)
%=-=-=-=-=-=-=-=-=-=-=-=-=-=-=-=-=-=-=-=-=-=-=-=-=-=-=-=-=-=-=-=-=-=-=-=-=-=-=-=-
% caption: Figure 1 -> <bold> Fig. 1 </bold>
\captionsetup[figure]{labelfont={bf},labelformat={default},labelsep=period,name={Fig.}}


%=-=-=-=-=-=-=-=-=-=-=-=-=-=-=-=-=-=-=-=-=-=-=-=-=-=-=-=-=-=-=-=-=-=-=-=-=-=-=-=-
% TITLE PAGE
%=-=-=-=-=-=-=-=-=-=-=-=-=-=-=-=-=-=-=-=-=-=-=-=-=-=-=-=-=-=-=-=-=-=-=-=-=-=-=-=-
%%%%%%%%%%%%%%%%%%%%%%%%%%%%%%%%%%%%%%%%%%
% Formal Book Title Page
% LaTeX Template
% Version 2.0 (23/7/17)
%
% This template was downloaded from:
% http://www.LaTeXTemplates.com
%
% Original author:
% Peter Wilson (herries.press@earthlink.net) with modifications by:
% Vel (vel@latextemplates.com)
%
% License:
% CC BY-NC-SA 3.0 (http://creativecommons.org/licenses/by-nc-sa/3.0/)
% 
% This template can be used in one of two ways:
%
% 1) Content can be added at the end of this file just before the \end{document}
% to use this title page as the starting point for your document.
%
% 2) Alternatively, if you already have a document which you wish to add this
% title page to, copy everything between the \begin{document} and
% \end{document} and paste it where you would like the title page in your
% document. You will then need to insert the packages and document 
% configurations into your document carefully making sure you are not loading
% the same package twice and that there are no clashes.
%
%%%%%%%%%%%%%%%%%%%%%%%%%%%%%%%%%%%%%%%%%

%----------------------------------------------------------------------------------------
%	PACKAGES AND OTHER DOCUMENT CONFIGURATIONS
%----------------------------------------------------------------------------------------



%----------------------------------------------------------------------------------------
%	TITLE PAGE
%----------------------------------------------------------------------------------------



\begin{titlepage} % Suppresses headers and footers on the title page

	\centering % Centre everything on the title page
	
	\scshape % Use small caps for all text on the title page
	
	\vspace*{\baselineskip} % White space at the top of the page
	
	%------------------------------------------------
	%	Title
	%------------------------------------------------
	
	\rule{\textwidth}{1.6pt}\vspace*{-\baselineskip}\vspace*{2pt} % Thick horizontal rule
	\rule{\textwidth}{0.4pt} % Thin horizontal rule
	
	\vspace{0.75\baselineskip} % Whitespace above the title
	
	{\LARGE COMPUTER VISION NOTES\\ \Large OBJECT LOCALISATION AND TRACKING\\} % Title
	
	\vspace{0.75\baselineskip} % Whitespace below the title
	
	\rule{\textwidth}{0.4pt}\vspace*{-\baselineskip}\vspace{3.2pt} % Thin horizontal rule
	\rule{\textwidth}{1.6pt} % Thick horizontal rule
	
	\vspace{2\baselineskip} % Whitespace after the title block
	
	%------------------------------------------------
	%	Subtitle
	%------------------------------------------------
	My personal notes on
	
	\vspace*{3\baselineskip} % Whitespace under the subtitle
	
	Object Localisation Techniques; Colour Matching, Mean Shift Tracking, Optical Flow, Lukas Kanade 
	
	\vspace*{3\baselineskip} % Whitespace under the subtitle
	
	%------------------------------------------------
	%	Editor(s)
	%------------------------------------------------
	
	By
	
	\vspace{0.5\baselineskip} % Whitespace before the editors
	
	{\normalfont \Large \mintinline{latex}{0xLeo} (\url{github.com/0xleo}) \\} % Editor list
	
	\vspace{0.5\baselineskip} % Whitespace below the editor list
	
	%\textit{The University of California \\ Berkeley} % Editor affiliation
	
	\vfill % Whitespace between editor names and publisher logo
	
	%------------------------------------------------
	%	Publisher
	%------------------------------------------------
	
	
	\vspace{0.3\baselineskip} % Whitespace under the publisher logo
	
	\today % Date
	
	{DRAFT X.YY} % Draft version
	{\\Missing: \ldots}

\end{titlepage}

%----------------------------------------------------------------------------------------

%\maketitle



%=-=-=-=-=-=-=-=-=-=-=-=-=-=-=-=-=-=-=-=-=-=-=-=-=-=-=-=-=-=-=-=-=-=-=-=-=-=-=-=-
% MAIN DOCUMENT
%=-=-=-=-=-=-=-=-=-=-=-=-=-=-=-=-=-=-=-=-=-=-=-=-=-=-=-=-=-=-=-=-=-=-=-=-=-=-=-=-




%------------------------------ New section ------------------------------%
\boxsimple{

Let 
\[
\bM = \begin{bmatrix}
1 & 0 & 1\\
1 & 1 & 0\\
0 & 1 & 1
\end{bmatrix}
\]
Find $\bM^{25}$  (\href{https://www.youtube.com/watch?v=7k2HmZJ_dYE}{source}).
}
\textbf{Solution:}

My solution is slightly different than the one in the video in the sense that it uses elementary matrix operations.

Matrix $\bM$ can be written as the sum of the unit matrix $\bI_3$ and another ``complement'' matrix $\bC$:
\begin{equation*}
    \bM = \bI_3 + \bC, \quad \bC = 
    \begin{bmatrix}
        0 & 0 & 1\\
        1 & 0 & 0\\
        0 & 1 & 0
    \end{bmatrix}
\end{equation*}

First, to expand $(\bI_3 + \bC)^{25}$, we use the binomial theorem (see \ref{app:binomial}):
\[
    (\bI + \bC)^{25} = \sum_{k=0}^{25}\frac{25!}{k!(25-k)!}\bI^{25-k}\bC^k = \sum_{k=0}^{25}\frac{25!}{k!(25-k)!}\bC^k
    \label{eq:sum_wrt_c}
    \tag{1}
\]
Now $\bC^k$ needs to be computed for each $k$. $\bC$ resembles with the $\bI_3$ unit matrix; we can transform $\bI_3$ to $\bC$ if we swap rows 1, 3 (operation $R_1 \leftrightarrow R_3$) and then swap rows 2, 3 (operation $R_2 \leftrightarrow R_3$):
\[
   \begin{bmatrix}
        1 & 0 & 0\\
        0 & 1 & 0\\
        0 & 0 & 1
    \end{bmatrix}
    \xrightarrow[]{R_1 \leftrightarrow R_3}
    \begin{bmatrix}
        0 & 0 & 1\\
        0 & 1 & 0\\
        1 & 0 & 0
    \end{bmatrix}
    \xrightarrow[]{R_2 \leftrightarrow R_3}
    \begin{bmatrix}
        0 & 0 & 1\\
        1 & 0 & 0\\
        0 & 1 & 0
    \end{bmatrix}
\]
Operation $R_{13}$ can be performed by matrix multiplication; by multiplying $\bI_3$ with a $3\times 3$ so called ``elementary'' matrix $\textbf{E}_{13}$ such that it swaps rows 1 and 3. So does $R_{23}$, but with $\textbf{E}_{23}$. Therefore, $\bC$ can be expressed with matrix multiplication as
\[
\bC = \textbf{E}_{23}\textbf{E}_{13} \bI_3 = \textbf{E}_{23}\textbf{E}_{13} 
\label{eq:mat_c}
\tag{2}
\]
One can find out (see \ref{app:elem_row}) that the sought row-swapping matrices $\textbf{E}_{23}$, $\textbf{E}_{13}$ are:
\[
    \textbf{E}_{23} = 
    \begin{bmatrix}
        1 & 0 & 0 \\
        0 & 0 & 1 \\
        0 & 1 & 0
    \end{bmatrix}, \quad
    \textbf{E}_{13} = 
    \begin{bmatrix}
        0 & 0 & 1 \\
        0 & 1 & 0 \\
        1 & 0 & 0
    \end{bmatrix}
\]
Computing the next higher power:
\begin{align*}
\bC^2 &= \textbf{E}_{23}\textbf{E}_{13}\bC =   \textbf{E}_{23}\textbf{E}_{13}\begin{bmatrix} 
        0 & 0 & 1\\
        1 & 0 & 0\\
        0 & 1 & 0
    \end{bmatrix} 
= \textbf{E}_{23}
    \begin{bmatrix}
        0 & 1 & 0\\
        1 & 0 & 0\\
         0 & 0 & 1\\
    \end{bmatrix}\\
&= 
    \begin{bmatrix}
        0 & 1 & 0\\
        0 & 0 & 1\\
        1 & 0 & 0\\
    \end{bmatrix}
    \label{eq:mat_c_squared}
    \tag{3}
\end{align*}
Raising to the next higher power, we obtain a nice recursion:
\begin{align*}
    \bC^3 &= \bC\bC^2\\
    &= \textbf{E}_{23}\textbf{E}_{13}\bC^2 = \textbf{E}_{23}\textbf{E}_{13} \begin{bmatrix}
            0 & 1 & 0\\
            0 & 0 & 1\\
            1 & 0 & 0\\
        \end{bmatrix}\\
    &=\textbf{E}_{23} \begin{bmatrix}
            1 & 0 & 0\\
            0 & 0 & 1\\
            0 & 1 & 0\\
        \end{bmatrix}
    = \begin{bmatrix}
            1 & 0 & 0\\
            0 & 1 & 0\\
            0 & 0 & 1\\
        \end{bmatrix} \\
    &= \bI_3
    \label{eq:mat_c_cube}
    \tag{4}
\end{align*}
In general, as derived from \eqref{eq:mat_c}, \eqref{eq:mat_c_squared}, \eqref{eq:mat_c_cube} powers of $\bC$ are repeated every 3:
\[ 
\bC^n = \left\{
\begin{array}{ll}
    \bI_3, & n = 3k, \quad k \in \setN\\
    \bC, & n = 3k+1, \quad k \in \setN\\
    \bC^2, & n = 3k+2, \quad k \in \setN\\
\end{array} 
\right. 
\]
Both $\bC$ and $\bC^2$ have been computed so the sum in \eqref{eq:sum_wrt_c} can be calculated in terms of known matrices. Notice that the term $\bC^n$ of the sum will equal to $\bI_3$ for $n=0,3,\ldots,24$, to $\bC$ for $n=1,4,\ldots,23$, and to $\bC^2$ for $n=2,5,\ldots,25$. Therefore we break the sum in 3 components, each one running for each indexing mentioned above:
\[
    \sum_{k=0}^{25}\frac{25!}{k!(25-k)!}\bC^k =
    \underbrace{\sum_{\substack{k=0, \\ k+=3}}^{24}\frac{25!}{k!(25-k)!}}_{\alpha_1}\bI_3 +
    \underbrace{\sum_{\substack{k=1, \\ k+=3}}^{23}\frac{25!}{k!(25-k)!}}_{\alpha_2}\bC +
    \underbrace{\sum_{\substack{k=2, \\ k+=3}}^{25}\frac{25!}{k!(25-k)!}}_{\alpha_3}\bC^2
\]
So we know in advance that the final will look like in terms of $\alpha_1,\; \alpha_2, \; \alpha_3$ as follows:
\[
     \sum_{k=0}^{25}\frac{25!}{k!(25-k)!}\bC^k = 
     \begin{bmatrix}
         \alpha_1 & \alpha_3 & \alpha_2\\
         \alpha_2 & \alpha_1 & \alpha_3 \\
         \alpha_3 & \alpha_2 & \alpha_1
     \end{bmatrix}
     \label{eq:sum_as_parameter_matrix}
     \tag{5}
\]
$\alpha_1,\; \alpha_2,\; \alpha_3$ can easily be computed e.g. in Python3:
\begin{verbatim}
from math import factorial as fac
a1 = sum([fac(25)/(fac(k)*fac(25-k)) for k in list(range(0,26,3))])
a2 = sum([fac(25)/(fac(k)*fac(25-k)) for k in list(range(1,26,3))])
a3 = sum([fac(25)/(fac(k)*fac(25-k)) for k in list(range(2,26,3))])
\end{verbatim}
We obtain that $\alpha_1 = \alpha_2 = 11184811$ and $\alpha_3 = 11184810$, therefore using \eqref{eq:sum_as_parameter_matrix} the final answer is
\begin{align*}
      \bM^{25} &= 
     \begin{bmatrix}
        1 & 0 & 1\\
        1 & 1 & 0\\
        0 & 1 & 1
    \end{bmatrix}^{25}  \\
     &=\sum_{k=0}^{25}\frac{25!}{k!(25-k)!}\bC^k = 
     \begin{bmatrix}
         11184811 & 11184810 & 11184811\\
         11184811 & 11184811 & 11184810 \\
         11184810 & 11184811 & 11184811
     \end{bmatrix}   
\end{align*}




%=-=-=-=-=-=-=-=-=-=-=-=-=-=-=-=-=-=-=-=-=-=-=-=-=-=-=-=-=-=-=-=-=-=-=-=-=-=-=-=-
% References
%=-=-=-=-=-=-=-=-=-=-=-=-=-=-=-=-=-=-=-=-=-=-=-=-=-=-=-=-=-=-=-=-=-=-=-=-=-=-=-=-
\newpage
\printbibliography





%=-=-=-=-=-=-=-=-=-=-=-=-=-=-=-=-=-=-=-=-=-=-=-=-=-=-=-=-=-=-=-=-=-=-=-=-=-=-=-=-
% Appendices
%=-=-=-=-=-=-=-=-=-=-=-=-=-=-=-=-=-=-=-=-=-=-=-=-=-=-=-=-=-=-=-=-=-=-=-=-=-=-=-=-
\newpage
\appendix

\section{Appendices}
% ------------------------ New appendix ------------------------ %
\newpage
\subsection{Binomial theorem}
\label{app:binomial}

\begin{theorem}[Binomial Theorem]
Let $x,y\in\setR$ and $n\geq 0$. Then
\begin{equation*}
   (x+y)^n = \sum_{k=0}^n\binom{n}{k}x^{n-k}y^k, \quad \binom{n}{k} = \frac{n!}{k!(n-k)!}
\end{equation*}
\end{theorem}

\begin{proof}
It is derived by induction as described in \cite{amsi}.
\end{proof}

% ------------------------ New appendix ------------------------ %
\newpage
\subsection{Elementary row interchange operations}
\label{app:elem_row}

We will observe and derive the elementary row interchange matrices $\textbf{E}_{ij}$, $i,j\leq 3, i \neq j$ for $3 \times 3$ matrices, however they can easily be generalised for $m \times n$. Supposed we are given a
\[
\bM = 
\begin{bmatrix}
a & b & c \\
d & e & f \\
g & h & i
\end{bmatrix}
\]
and we want to generate a new matrix $\textbf{N}$ whose 1st row is $\bM$'s 3rd row and is otherwise empty. Consider the row vectors $\be_1 = \begin{bmatrix}1 & 0 & 0\end{bmatrix}$, $\be_2 = \begin{bmatrix}0 & 1 & 0\end{bmatrix}$, $\be_3 = \begin{bmatrix}0 & 0 & 1\end{bmatrix}$. Multiplying $\be_1$ by $\textbf{M}$ extracts the 1st row of $\bM$, and so on for $\be_2$, $\be_3$:
\[
\be_1 \bM = 
\begin{bmatrix}
1 &0 &0
\end{bmatrix}
\begin{bmatrix}
a & b & c \\
d & e & f \\
g & h & i
\end{bmatrix} = 
\begin{bmatrix}
a & b & c
\end{bmatrix},
\quad \be_2 \bM = \begin{bmatrix} d & e & f \end{bmatrix}, \quad
\quad \be_3 \bM = \begin{bmatrix} g & h & i \end{bmatrix}, \quad
\]
Therefore if we multiply $\textbf{E}_{13} := \begin{bmatrix}
\be_3 & \be_2 & \be_2
\end{bmatrix}\t$ with $\bM$, the result will consist of $\bM$'s 3rd row  stacked on top of its 2nd row, stacked on top of the 1st row, i.e. $\bM$ with its rows 1 and 3 swapped:
\[
\textbf{E}_{13}\bM = 
\begin{bmatrix}
\textbf{e}_3 \\ \textbf{e}_2 \\ \textbf{e}_1
\end{bmatrix} \bM = 
\begin{bmatrix}
0 & 0 & 1\\
0 & 1 & 0 \\
1 & 0 & 0
\end{bmatrix}\bM = 
\begin{bmatrix}
g & h & i\\
d & e & f \\
a & b & c 
\end{bmatrix}
\]
Similarly, $\textbf{E}_{23}\bM$ swaps rows 2 and 3 of matrix $\bM$ and $\textbf{E}_{23}\textbf{E}_{13}\bM$ swaps row 1 with 3, and then rows 2 with 3:
\begin{gather*}
\textbf{E}_{23}\textbf{E}_{13}\bM = 
\textbf{E}_{23}
\begin{bmatrix}
g & h & i\\
d & e & f \\
a & b & c 
\end{bmatrix} = 
\begin{bmatrix}
g & h & i\\
a & b & c \\
d & e & f 
\end{bmatrix}
\end{gather*}

There's a neat property to note about the row interchange matrix $\textbf{E}_{ij}$; its inverse is itself. That should be intuitively obvious, as multiplying .e.g $\textbf{E}_{13}$ by $\bM$ swaps rows 1, 3 and multiplying again by $\textbf{E}_{13}$ swaps again rows 1,3, leading back to the original matrix $\bM$.
\begin{lemma}[inverse of row interchange matrix]
Let  $\textbf{E}_{ij}$ be an elementary row interchange matrix that swaps rows $i, j, i,j \geq m$ of an invertible matrix $\bM\in \setR^{m\times n}$. Then:
\begin{equation}
    \textbf{E}_{ij} = \textbf{E}_{ij}^{-1}
\end{equation}
\end{lemma}
\begin{proof}
\begin{gather*}
    \textbf{E}_{ij}\textbf{E}_{ij}\bM = \bM \Rightarrow \\
    \textbf{E}_{ij}\textbf{E}_{ij} = \bI_m
    \label{eq:elem_matrix_inv}
\end{gather*}
\end{proof}
Some other properties that can be readily derived from \eqref{eq:elem_matrix_inv} are:
% https://en.wikipedia.org/wiki/Exchange_matrix transpose and powers
\begin{lemma}[properties of row interchange matrices]
If $\textbf{E}_{ij}$ is an elementary matrix that interchanges rows $i$ and $j$, then:
\begin{equation}
    \textbf{E}_{ij}^n =
    \left\{
    \begin{array}{ll}
    \textbf{E}_{ij}, & n = 2k, \quad k \in \setN\\
    \bI , & n = 2k+1, \quad k \in \setN
\end{array} 
\right.
\end{equation}
Furthermore, $\textbf{E}_{ij}$ is symmetric:
\begin{equation}
    \textbf{E}_{ij}\t = \textbf{E}_{ij}
\end{equation}
\end{lemma}

\end{document}